\documentclass{article}
\usepackage{booktabs}
\usepackage{amssymb}
\usepackage[a4paper]{geometry}
\usepackage{fancyhdr}
\pagestyle{fancy}
\lhead{Simulationen}
\rhead{März 2026}
\begin{document}
\section{Simulationen}
Mit dem Taschenrechner und anderen Programmen kann ein Zufallsexperiment \emph{simuliert} werden.
 
\subsection{Werte Erstellen} 
Hierfür bietet der CAS zwei Funktionen:
\begin{center}
 \begin{tabular}{llll}
  Art & Funktion & Ergebnis & Menu \\
  \toprule
  binomial & \texttt{rand(a,b)} & Zahl $r \in [0;1] \in \mathbb{R}$ & 5, 4, 2\\
  normal & \texttt{rand()} & Zahl $r \in [a;b] \in \mathbb{N}$ 5, 4, 1 \\
  \bottomrule
 \end{tabular}
\end{center} 
Beide dieser Funktionen kann um ein Paramter $n$ erweitert werden, wird dieser genutzt, so gibt \texttt{rand(n)} beziehungsweise \texttt{randInt(a,b,n)} jeweils $n$ dieser zufälligen Zahlen an.
 
\subsection{Werte verarbeiten}
Mit \texttt{sum(list)} kann die Summe der Werte eine Liste gefunden werden.
 
Viel intressanter kann gezählt werden, wie viele der Elemente eine bestimmte Bedingung erfüllen, mit \texttt{countIf(liste,bedingung)}. Die Bedingung ist dabei in der Regel eine normale Gleichung, der zu überprüfende Wert der Liste wird dabei mit einem \texttt{?} gekennzeichnet.
\end{document}